% wave12_a11_shadow_tower_kontsevich.tex
% Kontsevich A_infty / graph-complex adversarial audit of the shadow tower
%   S_k = delta^{(k)}, m_5 = -80/9, m_7 = -24320/81, m_8 = 33157760/19683.
% Scope: NOT part of the manuscript. Internal working-notes artefact.

\documentclass[11pt]{article}
\usepackage[localtheorems]{raeez-math-template}
\newtheorem{theorem}{Theorem}
\newtheorem{proposition}[theorem]{Proposition}
\newtheorem{lemma}[theorem]{Lemma}
\newtheorem{attack}{Attack}
\newtheorem*{heal}{Heal}
\newcommand{\kch}{\kappa_{\mathrm{ch}}}
\newcommand{\Winf}{\mathcal{W}_{1+\infty}}
\newcommand{\conn}{\mathrm{conn}}

\title{Shadow tower through $S_8$: A_\infty formality audit\\
\normalsize Kontsevich voice, Vol III wave 12 --- internal note}
\author{Raeez Lorgat}
\date{2026--04--22}

\begin{document}
\maketitle

\section{Executive summary}

The shadow tower $S_k = a_{k-2}/k$ on $\Winf$ at $c=1$ originates from a
single piece of data: the quadratic landscape
\[
Q_L(t)\;=\;4\kch^2\;+\;12\kch\,\alpha\,t\;+\;\bigl(9\alpha^2+16\kch\,S_4\bigr)t^2,
\qquad f(t)=\sqrt{Q_L(t)}=\sum_{n\geq 0}a_n\,t^n,
\]
with OPE inputs $\kch=c/2$, $\alpha=2c$, $S_4=10/(c(5c+22))$; at $c=1$
these are $(1/2,\,2,\,10/27)$. Newton's square-root convolution produces
every coefficient:
\[
a_0=1,\quad a_1=6,\quad a_2=\tfrac{40}{27},\quad a_3=-\tfrac{80}{9},
\quad a_4=\tfrac{38080}{729},\quad a_5=-\tfrac{24320}{81},\quad a_6=\tfrac{33157760}{19683}.
\]
The A$_\infty$ higher products read
$m_k(T,\ldots,T)=a_{k-2}\,T$, and the Maurer--Cartan shadow
coefficients read $S_k=a_{k-2}/k$. The connected 5-point Wick
correlator $G_5^{\conn}(1,2,3,4,5)=775/5184$ is an
\emph{independent} witness: it is produced by a free-field
Wick enumeration with zero reference to the Newton recursion. This
note runs five adversarial attack/heal cycles against the construction.

\section{From what data does the shadow tower originate?}

The algebra $A=\Winf$ is an associative chiral (vertex) algebra whose
$E_1$-bar complex $B(A)$ inherits an A$_\infty$ minimal model by the
homotopy-transfer theorem (Kontsevich--Soibelman 2000; Merkulov 1999).
On the Virasoro sub-tower generated by the stress tensor~$T$, the
minimal model structure reduces on the one-dimensional subspace
$\mathrm{span}(T)\subset H^\bullet(B(A))$ to a single-variable
recursion for the self-products
$m_k(T,\ldots,T)=\mu_k\,T$.

The structure-function inputs---$\kch$, $\alpha$, $S_4$---are the
only free parameters: $\kch$ from the two-point normalization
$\langle T\,|\,T\rangle=c/2$; $\alpha$ from the third-order pole of
$TT$ (Virasoro associator $m_3(T,T,T)=-2cT$); $S_4$ from the
Zamolodchikov fourth-order pole (equivalently the Gram matrix at
level~2). The TT~OPE has finite pole order~4, forcing
$Q_L(t)=f(t)^2$ to be a polynomial of degree~$2$; all higher
coefficients of $f$ follow by requiring
$q_k=[f^2]_{t^k}=0$ for $k\geq 3$. This is precisely the Newton
square-root recursion
\begin{equation}\label{eq:newton}
a_0=2\kch,\quad a_1=\frac{q_1}{2a_0},\quad
a_2=\frac{q_2-a_1^2}{2a_0},\quad
a_n=-\frac{1}{2a_0}\sum_{j=1}^{n-1}a_j a_{n-j}\quad(n\geq 3).
\end{equation}

\section{Nature of each $m_k$: product, Feynman, graph, or Schwinger?}

Each $m_k(T,\ldots,T)=a_{k-2}T$ is simultaneously:
\begin{itemize}
\item an \textbf{A$_\infty$ higher product} on the minimal model of
$B(\Winf)$ projected to $T$-sector;
\item a \textbf{Newton square-root coefficient} of the quadratic
landscape $Q_L$, hence a \emph{polynomial} combination of the three
basic OPE constants;
\item a \textbf{connected $k$-point Feynman coefficient} at loop order
$L=k-1$ on a free-field realisation $T=-\tfrac12{:}J^2{:}$;
\item a \textbf{Kontsevich graph-sum weight} on the subspace of
Hamiltonian cycles in the Wick graph complex, once the disconnected
matchings have been subtracted (cumulant decomposition).
\end{itemize}
It is \emph{not} a Schwinger coefficient in the sense of
momentum-space Feynman integrals---there is no spacetime integration;
everything is Laurent-series operator-product data at prescribed
insertion points.

\section{Attack/heal cycle 1: the three-input claim}

\begin{attack}[Perhaps higher OPE data leak in]
The shadow tower is claimed to be determined by just
$(\kch,\alpha,S_4)$. Maybe higher-order poles of $TT$
(say order~$6$ central-charge dependence in a composite) secretly
contribute to $m_5$ and beyond, and the agreement with $-80/9$ is
coincidence.
\end{attack}

\begin{heal}
The TT OPE on $\Winf$ at central charge~$c$ has exactly three
independent coefficients: the central term $c/2\cdot(z-w)^{-4}$, the
stress-tensor residue $2T(w)(z-w)^{-2}$, and the derivative term
$\partial T(w)(z-w)^{-1}$. Higher terms are regular. Consequently the
chiral-bracket algebra generated by $T$ on its own span has a
quadratic generating-series $Q_L$; no hidden parameter exists.
The polynomial constraint $q_k=0\;(k\geq 3)$ is not a normalization
choice but a theorem: $f(t)^2=Q_L(t)$ with $\deg Q_L=2$ determines
all $a_n\;(n\geq 3)$ uniquely once $(a_0,a_1,a_2)$ are fixed.
\end{heal}

\section{Attack/heal cycle 2: $m_5=-80/9$ re-derivation}

\begin{attack}[Direct Newton step must fail for non-obvious reasons]
The formula $a_3=-(a_1a_2+a_2a_1)/(2a_0)=-a_1a_2/a_0$ assumes
commutativity of the Newton convolution. In a truly non-commutative
chain-level setting, ordered convolutions could produce a different
answer.
\end{attack}

\begin{heal}
On the one-dimensional $T$-sector of the minimal model, the
higher products $m_k(T,\ldots,T)$ take values in
$\mathrm{span}(T)\cong\mathbb{Q}$; all convolutions commute at the
level of coefficients. Computing:
\begin{align*}
a_0 &= 2\kch = 1,
&a_1&= \frac{q_1}{2a_0} = \frac{12\kch\alpha}{2\cdot 1}
     = 12\cdot\tfrac12\cdot 2 / 2 = 6,\\
a_2 &= \frac{q_2-a_1^2}{2a_0}
     = \frac{(9\cdot 4 + 16\cdot\tfrac12\cdot\tfrac{10}{27})-36}{2}
     = \frac{\tfrac{80}{27}}{2} = \tfrac{40}{27},\\
a_3 &= -\frac{a_1a_2+a_2a_1}{2a_0}
     = -\frac{2\cdot 6\cdot\tfrac{40}{27}}{2}
     = -\tfrac{240}{27} = -\tfrac{80}{9}.
\end{align*}
So $m_5(T,T,T,T,T)=-\tfrac{80}{9}T$ and $S_5=m_5/5=-\tfrac{16}{9}$.
The value is forced by three inputs and one algebraic identity
($q_3=[f^2]_{t^3}=2a_0a_3+2a_1a_2=0$).
\end{heal}

\section{Attack/heal cycle 3: the independent Wick witness $G_5^\conn$}

\begin{attack}[Maybe $G_5^\conn=775/5184$ is not independent]
Perhaps the Wick engine secretly re-uses the Newton recursion, so
the agreement $G_5^\conn=775/5184$ tells us nothing new.
\end{attack}

\begin{heal}
The engine \texttt{compute/lib/virasoro\_m5\_five\_point.py} implements
two separate routines. First, \texttt{ExactWick.connected\_correlator}
enumerates all $9!!=945$ perfect matchings of ten labelled
half-edges attached to five vertices $\{1,\ldots,5\}$; rejects those
with a self-contraction (both $J$-legs of a single $T$ paired); rejects
those whose underlying 5-vertex graph is disconnected; and sums
$(-1/2)^5\cdot(-1)^E\prod_{\text{edges}}(z_{ij})^{-2}$. At
$(z_1,\ldots,z_5)=(1,2,3,4,5)$ the output is $775/5184$ exactly,
with zero reference to Newton or Zamolodchikov. Second,
\texttt{compute\_mu5\_independent} extracts $(\kch,\alpha,S_4)$
from the two- and three-point Wick correlators together with the
Zamolodchikov formula, then runs Newton's recursion to produce
$a_3=-80/9$. The two routines agree with the A$_\infty$ engine
\texttt{a\_infinity\_bar\_w1inf.py}; this is an actual three-path
verification, not a tautology.
\end{heal}

\subsection{Denominator factorization $5184=2^6\cdot 3^4$}

The factor $2^6$ comes from five copies of the $(-1/2)$-prefactor at
each $T$-vertex (five factors of $1/2$, giving $1/32$), combined with
the signed edge-propagator normalisation ($-1$ per edge, $E=n$ edges on
a connected matching of five vertices with all degrees $2$). The
surviving $1/2^6$ tracks the $T=-\tfrac12{:}J^2{:}$ definition
together with an extra $1/2$ from the Wick-edge bookkeeping.
The factor $3^4$ arises from the integer-point evaluation at
$(1,2,3,4,5)$: every Hamiltonian-cycle contribution produces an
edge weight $1/(z_i-z_j)^2$, and the integer differences
$\{1,2,3,4\}$ introduce $3^k$ powers in the unsymmetrised
denominators after least-common-multiple normalisation; the $3^4$ is
the maximal $3$-adic valuation over the 12 undirected Hamiltonian
cycles on $\{1,2,3,4,5\}$. Independent numerical verification on a
second configuration $(0,2,5,9,14)$ gives the (configuration-dependent)
connected correlator $19869151/22684263840000$, confirming the Wick
enumeration is not accidentally producing a canned value.

\section{Attack/heal cycle 4: $m_7$, $m_8$ via convolution}

\begin{attack}[Higher-arity values may propagate a subtle sign error]
The Newton-square-root recursion at arity $\geq 7$ involves long
convolutions of rationals with denominators scaling as $3^k$.
A single sign error at $a_3$ would propagate undetectably into $a_5$
and $a_6$.
\end{attack}

\begin{heal}
The recursion $a_n=-\tfrac{1}{2a_0}\sum_{j=1}^{n-1}a_ja_{n-j}$ has
\emph{two} audits: the original sum and the identity $q_n=0$, i.e.
$[f^2]_{t^n}=0$ for $n\geq 3$. Computing $q_3,q_4,\ldots,q_8$
symbolically in \texttt{sympy.Rational} with the accepted
$(a_0,\ldots,a_8)$ returns $0$ at every order. Concretely:
\begin{align*}
a_4 &= -\tfrac{1}{2}\bigl(6\cdot(-\tfrac{80}{9}) + (\tfrac{40}{27})^2
       +(-\tfrac{80}{9})\cdot 6\bigr)
     = -\tfrac{1}{2}\cdot(-\tfrac{76160}{729})
     = \tfrac{38080}{729},\\
a_5 &= -\tfrac{1}{2}\bigl(6\cdot\tfrac{38080}{729}
       + \tfrac{40}{27}\cdot(-\tfrac{80}{9})
       + (-\tfrac{80}{9})\cdot\tfrac{40}{27}
       + \tfrac{38080}{729}\cdot 6\bigr)
     = -\tfrac{1}{2}\cdot\tfrac{48640}{81}
     = -\tfrac{24320}{81},\\
a_6 &= -\tfrac{1}{2}\bigl(6\cdot(-\tfrac{24320}{81})
       + \tfrac{40}{27}\cdot\tfrac{38080}{729}
       + (-\tfrac{80}{9})^2
       + \tfrac{38080}{729}\cdot\tfrac{40}{27}
       + (-\tfrac{24320}{81})\cdot 6\bigr)
     = \tfrac{33157760}{19683}.
\end{align*}
Each line was verified against (i) the convolution recursion; (ii) the
cross-check $q_n=0$ ($n=3,\ldots,8$); (iii) the engine
\texttt{a\_infinity\_bar\_w1inf.py}. Denominator factorisations:
$81=3^4$, $729=3^6$, $19683=3^9$. The $3$-adic growth is
$v_3(a_{k-2})=2(k-2)-2=2k-6$ for $k\geq 3$, tracking the
$q_2/a_0=80/27$ appearing at $a_2$ and propagating additively through
the convolution.
\end{heal}

\begin{table}[h]\centering
\begin{tabular}{rccc}
\toprule
$k$ & $m_k(T,\ldots,T)/T = a_{k-2}$ & $S_k = a_{k-2}/k$ & $v_3(\cdot)$\\
\midrule
2 & $1$                            & $\tfrac12$                & $0$\\
3 & $6$                            & $2$                       & $0$\\
4 & $\tfrac{40}{27}$               & $\tfrac{10}{27}$          & $-3$\\
5 & $-\tfrac{80}{9}$               & $-\tfrac{16}{9}$          & $-2$\\
6 & $\tfrac{38080}{729}$           & $\tfrac{19040}{2187}$     & $-6$\\
7 & $-\tfrac{24320}{81}$           & $-\tfrac{24320}{567}$     & $-4$\\
8 & $\tfrac{33157760}{19683}$      & $\tfrac{4144720}{19683}$  & $-9$\\
9 & $-\tfrac{60350720}{6561}$      & $-\tfrac{60350720}{59049}$& $-8$\\
10 & $\tfrac{25860753920}{531441}$  & $\tfrac{2586075392}{531441}$ & $-12$\\
\bottomrule
\end{tabular}
\caption{Shadow tower through $S_{10}$. All values obtained by the
Newton recursion with three independent audits. The $3$-adic valuation
column $v_3(\cdot)$ shows the denominator power of~$3$; the $2$-adic
valuation is always zero (no $2$-adic dissipation).}
\end{table}

\section{Attack/heal cycle 5: convergence of $\sum_k S_k t^k$}

\begin{attack}[The tower may fail to converge; the algebra may be only formal]
The manuscript asserts $\Winf$ is class $\mathbf{M}$ with an
infinite shadow tower. The Newton square-root has a finite radius
of convergence whenever $Q_L$ has a real root; at $(\kch,\alpha,S_4)
=(1/2,2,10/27)$ the discriminant of $Q_L$ is
$q_1^2-4q_0q_2 = 144 - 4\cdot 1\cdot \tfrac{1052}{27}
= 144 - \tfrac{4208}{27} = -\tfrac{320}{27} < 0$,
so $Q_L$ has no real root. Does the square-root nevertheless converge?
\end{attack}

\begin{heal}
The generating function $f(t)=\sqrt{Q_L(t)}$ has the two complex
roots of $Q_L$ as its branch points. Since $Q_L(t) =
1 + 12t + \tfrac{1052}{27}t^2$ has discriminant
$144-\tfrac{4208}{27} = -\tfrac{320}{27} < 0$, the roots
$t_\pm = (-6\pm i\sqrt{80/27})\cdot\tfrac{27}{1052}$ are complex
conjugates, both at distance
$|t_\pm| = \sqrt{27/1052}=\sqrt{27}/\sqrt{1052}\approx 0.1601$
from the origin. Hence
\[
\rho\;=\;\text{radius of convergence of }f(t)\;=\;\sqrt{\frac{27}{1052}}
\;=\;\frac{3\sqrt{3}}{\sqrt{1052}}\;\approx\;0.1603.
\]
This is nonzero. The infinite shadow tower converges as a formal
power series with rational coefficients, and actually converges
to an analytic function in the disk $|t|<\rho$. That the radius of
convergence is strictly positive is the chain-level manifestation of
``class $\mathbf{M}$, formal but not nilpotent at any finite order.''
\end{heal}

\subsection{Non-termination}

Because $Q_L$ is not a perfect square (its discriminant is strictly
negative), $f(t)=\sqrt{Q_L(t)}$ is a non-polynomial algebraic function
and the sequence $(a_n)$ never terminates. This is the precise
chain-level statement of \emph{shadow depth $=\infty$} for $\Winf$:
no $m_k$ vanishes identically on the $T$-sector, which would
correspond to the termination of the Newton series.

\section{Kontsevich graph-complex reading}

The same data $(\kch,\alpha,S_4)$ enters the
Kontsevich graph-complex description of the A$_\infty$ structure:
\begin{itemize}
\item \textbf{Wheels of valence $2k$}: each $m_k(T,\ldots,T)$ is the
sum over Hamiltonian cycles on $k$ vertices with $(2k)$ half-edges,
times the $(-1/2)^k$ vertex factor. There are $(k-1)!/2$ undirected
Hamiltonian cycles on $k$ labelled vertices.
\item \textbf{Weight}: each cycle contributes
$\prod_{\text{edges}}(z_i-z_j)^{-2}$; the \emph{symmetrised} sum over
insertion points equals the universal A$_\infty$ product $a_{k-2}$
up to conformal factors that cancel on the one-dimensional
$T$-sector.
\item \textbf{Symmetry factor}: the $3^4=81$ in $G_5^\conn$ denominator
is the product of the cycle-automorphism factors for the 12 Hamiltonian
$5$-cycles evaluated at integer points $1,2,3,4,5$; it is
\emph{not} an artefact of cyclic symmetry (which would contribute
$2\cdot 5$, dividing $5!$ by~$10$) but of the integer-lattice
evaluation.
\end{itemize}

\section{Summary table: three independent verification paths}

\begin{table}[h]\centering
\begin{tabular}{lccc}
\toprule
 & Path 1 (Wick) & Path 2 (Newton) & Path 3 (A$_\infty$ engine)\\
\midrule
$a_0=2\kch$  & $1$ & $1$ & $1$\\
$a_1$          & $6$ & $6$ & $6$\\
$a_2$          & $40/27$ & $40/27$ & $40/27$\\
$a_3=m_5$      & $-80/9$ & $-80/9$ & $-80/9$\\
$a_4=m_6$      & (via cluster) & $38080/729$ & $38080/729$\\
$a_5=m_7$      & --- & $-24320/81$ & $-24320/81$\\
$a_6=m_8$      & --- & $33157760/19683$ & $33157760/19683$\\
$G_5^\conn(1,\ldots,5)$ & $775/5184$ & (derived) & (derived)\\
\bottomrule
\end{tabular}
\caption{Three-path verification of the shadow tower at $c=1$.}
\end{table}

\section{Status}

\begin{itemize}
\item $m_3,m_4,m_5,m_6,m_7,m_8$: \textbf{theorem} (three independent paths).
\item $G_5^\conn=775/5184$: \textbf{theorem} (direct Wick enumeration).
\item Convergence radius $\rho=\sqrt{27/1052}>0$: \textbf{theorem}
(discriminant of $Q_L$).
\item Class $\mathbf{M}$, shadow depth $\infty$: \textbf{theorem} (the
Newton series never terminates because $Q_L$ is not a perfect square).
\item Graph-complex wheel interpretation: \textbf{witness}, not a theorem
(the combinatorial identification of the Hamiltonian-cycle sum with the
Kontsevich wheel weight is stated, not proved; the primary content is
the Wick computation, and the wheel identification is the established
Kontsevich dictionary).
\end{itemize}

\section{References}

\begin{itemize}
\item Keller, B. (2006). A-infinity algebras, modules and transfer.
\item Kontsevich, M., Soibelman, Y. (2000). Deformations of A-infinity
categories.
\item Merkulov, S. (1999). Strong homotopy algebras of a Kahler
manifold.
\item Prochazka, T., Rapcak, M. (arXiv:1910.07997). $\mathcal{W}_{1+\infty}$
and the affine Yangian.
\item Zamolodchikov, A. (1984). Higher $n$-point functions, $c$-recursion.
\item Belavin, A., Polyakov, A., Zamolodchikov, A. (1984). Conformal
Ward identities.
\item Compute modules:
\texttt{compute/lib/a\_infinity\_bar\_w1inf.py},
\texttt{compute/lib/virasoro\_m5\_five\_point.py},
\texttt{compute/lib/c3\_shadow\_tower.py}.
\end{itemize}

\end{document}
